% GENERATED FILE. Edit canonical lessons, then run npm run generate:books.
% canonical-source-hash: fnv1a64:df88e6de96a940b0
% canonical-lessons: FR-C296-braquage, FR-C296-braqueur, FR-C296-otage, FR-C296-rancon, FR-C296-echapper

\chapter{Le braquage}
\label{ch:fr-le-braquage}
\hlchaptermodality{296}

\begin{chapteropening}
\textbf{By the end of this chapter:} \emph{I can tell the story of a hold-up: the armed robbery, the robbers, a hostage, a ransom and an escape.}

Complete the last lesson of chapter 296: I can tell the story of a hold-up: the armed robbery, the robbers, a hostage, a ransom and an escape.
\end{chapteropening}

\section[le braquage]{le braquage — the armed robbery, the hold-up}

\label{lesson:FR-C296-braquage}



\begin{quote}
Before the new one: say the French for to acquit, then the French for to release.
\end{quote}

\subsection*{You'll want to know: le braquage}
\textbf{le braquage} — \textquotedblleft{}the armed robbery, the hold-up\textquotedblright{}.

\begin{quote}
Le braquage de la banque a duré dix minutes.
\end{quote}

\begin{quote}
\emph{The bank robbery lasted ten minutes.}
\end{quote}

Informal but frequent in the news.

\begin{cousinweb}
\emph{braquage} is built with \emph{-age} on the verb meaning \textquotedblleft{}to aim, to point (a gun)\textquotedblright{}; the origin of that verb is not settled (EtymDB lists no parent).
\end{cousinweb}

\begin{grammarlens}[title={What you've built}]
One more word for telling, explaining and describing at B1.
\end{grammarlens}

\subsection*{Your turn}
\begin{itemize}
\raggedright
  \item \emph{Say it:} \emph{le braquage}
  \item \emph{Say it:} \emph{le braquage}, once more
  \item \emph{Say it:} a sentence of your own with \emph{le braquage}
  \item \emph{Recall:} say the French for to acquit, then the French for to release, then say \emph{le braquage} again
\end{itemize}

\begin{tcolorbox}[breakable,colback=teal!4,colframe=teal!35!black,title={Before you move on}]
What does le braquage mean? (\textquotedblleft{}The armed robbery, the hold-up\textquotedblright{}.) Say it once more.
\end{tcolorbox}
\section[le braqueur]{le braqueur — the armed robber}

\label{lesson:FR-C296-braqueur}



\begin{quote}
Before the new one: say the French for to release, then the French for the armed robbery.
\end{quote}

\subsection*{You'll want to know: le braqueur}
\textbf{le braqueur} — \textquotedblleft{}the armed robber\textquotedblright{}.

\begin{quote}
Les braqueurs sont partis avec l'argent.
\end{quote}

\begin{quote}
\emph{The robbers left with the money.}
\end{quote}

Feminine: \emph{la braqueuse}.

\begin{cousinweb}
\emph{braqueur} is built with the agent suffix \emph{-eur} on the verb meaning \textquotedblleft{}to aim a gun\textquotedblright{}; the origin is not settled (EtymDB lists no parent).
\end{cousinweb}

\begin{grammarlens}[title={What you've built}]
One more word for telling, explaining and describing at B1.
\end{grammarlens}

\subsection*{Your turn}
\begin{itemize}
\raggedright
  \item \emph{Say it:} \emph{le braqueur}
  \item \emph{Say it:} \emph{le braqueur}, once more
  \item \emph{Say it:} a sentence of your own with \emph{le braqueur}
  \item \emph{Recall:} say the French for to release, then the French for the armed robbery, then say \emph{le braqueur} again
\end{itemize}

\begin{tcolorbox}[breakable,colback=teal!4,colframe=teal!35!black,title={Before you move on}]
What does le braqueur mean? (\textquotedblleft{}The armed robber\textquotedblright{}.) Say it once more.
\end{tcolorbox}
\section[l'otage]{l'otage — the hostage}

\label{lesson:FR-C296-otage}



\begin{quote}
Before the new one: say the French for the armed robbery, then the French for the armed robber.
\end{quote}

\subsection*{You'll want to know: l'otage}
\textbf{l'otage} — \textquotedblleft{}the hostage\textquotedblright{}.

\begin{quote}
Les braqueurs ont pris trois otages.
\end{quote}

\begin{quote}
\emph{The robbers took three hostages.}
\end{quote}

Masculine even when the hostage is a woman: \emph{un otage}.

\begin{cousinweb}
\emph{otage} comes from Old French \emph{ostage}, usually traced to \emph{oste} \textquotedblleft{}guest, host\textquotedblright{}, from Latin \emph{hospes} (Lewis \& Short): a guest held as a pledge. EtymDB lists no parent.
\end{cousinweb}

\begin{grammarlens}[title={What you've built}]
One more word for telling, explaining and describing at B1.
\end{grammarlens}

\subsection*{Your turn}
\begin{itemize}
\raggedright
  \item \emph{Say it:} \emph{l'otage}
  \item \emph{Say it:} \emph{l'otage}, once more
  \item \emph{Say it:} a sentence of your own with \emph{l'otage}
  \item \emph{Recall:} say the French for the armed robbery, then the French for the armed robber, then say \emph{l'otage} again
\end{itemize}

\begin{tcolorbox}[breakable,colback=teal!4,colframe=teal!35!black,title={Before you move on}]
What does l'otage mean? (\textquotedblleft{}The hostage\textquotedblright{}.) Say it once more.
\end{tcolorbox}
\section[la rançon]{la rançon — the ransom}

\label{lesson:FR-C296-rancon}



\begin{quote}
Before the new one: say the French for the armed robber, then the French for the hostage.
\end{quote}

\subsection*{You'll want to know: la rançon}
\textbf{la rançon} — \textquotedblleft{}the ransom\textquotedblright{}.

\begin{quote}
La famille a payé une rançon.
\end{quote}

\begin{quote}
\emph{The family paid a ransom.}
\end{quote}

\emph{La rançon du succès} — the price of success.

\begin{cousinweb}
\emph{rançon} comes from Old French \emph{rençon} (EtymDB), which is inherited from Latin \emph{redēmptiō} \textquotedblleft{}a buying back\textquotedblright{}, from \emph{redimere} (Lewis \& Short).
\end{cousinweb}

\begin{grammarlens}[title={What you've built}]
One more word for telling, explaining and describing at B1.
\end{grammarlens}

\subsection*{Your turn}
\begin{itemize}
\raggedright
  \item \emph{Say it:} \emph{la rançon}
  \item \emph{Say it:} \emph{la rançon}, once more
  \item \emph{Say it:} a sentence of your own with \emph{la rançon}
  \item \emph{Recall:} say the French for the armed robber, then the French for the hostage, then say \emph{la rançon} again
\end{itemize}

\begin{tcolorbox}[breakable,colback=teal!4,colframe=teal!35!black,title={Before you move on}]
What does la rançon mean? (\textquotedblleft{}The ransom\textquotedblright{}.) Say it once more.
\end{tcolorbox}
\section[échapper]{échapper — to escape; to slip from (one's hands)}

\label{lesson:FR-C296-echapper}



\begin{quote}
Before the new one: say the French for the hostage, then the French for the ransom.
\end{quote}

\subsection*{You'll want to know: échapper}
\textbf{échapper} — \textquotedblleft{}to escape; to slip from (one's hands)\textquotedblright{}.

\begin{quote}
Un des braqueurs s'est échappé par le toit.
\end{quote}

\begin{quote}
\emph{One of the robbers escaped over the roof.}
\end{quote}

Used both ways. With \emph{se} for breaking out: \emph{s'échapper de prison}. Without \emph{se}, with \emph{à}: \emph{échapper à un danger} — to escape a danger.

\begin{cousinweb}
\emph{échapper} comes from Old French \emph{eschaper} (EtymDB), usually traced to Vulgar Latin \emph{excappāre} \textquotedblleft{}to slip out of one's cloak\textquotedblright{}, built on \emph{cappa} \textquotedblleft{}a cloak\textquotedblright{}.
\end{cousinweb}

\begin{grammarlens}[title={What you've built}]
One more word for telling, explaining and describing at B1.
\end{grammarlens}

\subsection*{Your turn}
\begin{itemize}
\raggedright
  \item \emph{Say it:} \emph{échapper}
  \item \emph{Say it:} \emph{échapper}, once more
  \item \emph{Say it:} a sentence of your own with \emph{échapper}
  \item \emph{Recall:} say the French for the hostage, then the French for the ransom, then say \emph{échapper} again
\end{itemize}

\begin{tcolorbox}[breakable,colback=teal!4,colframe=teal!35!black,title={Before you move on}]
What does échapper mean? (\textquotedblleft{}To escape; to slip from (one's hands)\textquotedblright{}.) Say it once more.
\end{tcolorbox}
